Guo and Xu, Theorems 3.2--3.3, prove marginal and joint IPW-Z Gaussian limits after scaled inverse propensities stabilize; their Section 3.1 no-intercept clipped-LinUCB example supplies numerical evidence of two modes, not a proved finite-attractor theorem or certification of the present class. Guo--Xu is an informed-by comparator. Saco, arXiv:2602.15559v1, Theorem 5.14, asserts realized-QV and feasible Studentized normal limits and Wald coverage under Assumptions 3.6, 3.2, 3.7, 3.8, 4.5, 4.8, 3.9, 3.10, and 5.11, together with maintained Assumption 3.5, without deterministic variance stabilization. The bounded AIPW counterexample proposition satisfies those published assumptions but has a noncentered self-normalized limit, so Saco's theorem is a refuted comparator rather than an established valid pivot or a strict-tightening predecessor. The early-scale theorem supplies a sufficient early-measurability boundary and proves uniform validity on its stated class; no strict-tightening claim is made. The Abstract and Introduction of Shen and coauthors distinguish directional stability from full-matrix stability; that distinction informs Theorems 3--4, while the finite-basin every-level converses and bounded separation witness are not presented as a strict comparator chain. Han studies an orthogonal instability mechanism and target: for generalized Gaussian Thompson sampling, Proposition 2.6 and Theorem 2.7 on PDF pages 9--10, together with introduction equation (1.4) on page 4, identify the joint optimal-arm pull-proportion and noise limit with the unique invariant law of a continuous-state SDE. Han's Theorem 3.2 on PDF page 12 and introduction equation (1.5) on page 5 give Gaussian normalized-mean limits for stable arms and an SDE-defined semi-universal limit for multiple optimal arms; the body reports simulation evidence that the latter law is genuinely non-Gaussian, and Corollary 3.3 on PDF page 13 constructs confidence intervals from its quantiles. Han is close motivation, not a result anticipating Theorems 1--4 or the bounded Guo--Xu certification question. Leiner, Dunn, and Ramdas study off-policy M-estimation under working-model misspecification for adaptively collected bandit data. Their target is a projected solution under a stationary evaluation policy, and flexible variance stabilization permits unstable or nonconvergent logging policies; this is an informed-by comparator rather than a same-path-QV or finite-attractor result. Borkar supplies local lock-in bounds, and Pemantle supplies unstable-point avoidance, but neither certifies the bounded recursion. Classical Aldous--Eagleson stable martingale theory supplies the stable and mixed-normal limit apparatus. Theorem 1 instead uniformly derives the random basin covariance from finite separated attractors and a common Cesaro modulus, constructs an early \(\mathcal F_{k_T}\)-measurable decoder that avoids conditioning on terminal \(J\), and carries the limit through uniform IPW--Z exact-root linearization. Mourrat supplies only the bounded-martingale central-limit and quadratic-variation ingredient. Teicher together with Andrews and Mallows supplies identification of centered normal scale mixtures. Qian and coauthors and Liao and coauthors define the HeartSteps causal-excursion consumer. Zhang, Janson, and Murphy study pooled longitudinal adaptive sampling with Z-estimation and an adaptive sandwich estimator in a fixed-within-user, growing-participant regime. Thus the comparison supports a field-tier conditional theorem bundle without claiming a valid strict-tightening comparator.